\chapter{Sets and the operations that create them}\label{ch:sets}
\begin{readergoal}Use membership, inclusion, unions, intersections, set differences, power sets, and the idea of an admissible collection.\end{readergoal}
\section{From a collection to a set}
A set is a mathematical collection whose members are determined by a membership rule. If $U=\{a,b,c\}$, then $a\in U$ and $d\notin U$. In the ordinary language of mathematics, $a\in U$ means that $a$ \emph{belongs to} $U$; $A\subseteq U$ means every element of $A$ belongs to $U$. These statements differ: $\{a\}\subseteq U$, but typically $\{a\}\notin U$. The empty set $\varnothing$ has no members; $\{\varnothing\}$ has one member, which happens to be the empty set.

Two sets are equal exactly when they have identical members (the extensional criterion). Mathematical set construction must take place within an explicit foundational framework; one cannot freely declare every grammatical property to specify a set of all possible objects. The unrestricted collection ``all sets that are not members of themselves'' generates Russell's paradox. Following the accessible axiomatic spirit of \cite{HalmosNaive,MendelsonNumbers}, we allow subsets defined by conditions on an already established set and use conventional constructions for other sets.

\section{Set operations from ordinary English}
Let $A$ be the trials in which the animal runs, and $B$ those with a detected theta peak. Then $A\cap B$ consists of trials satisfying \emph{both} criteria, while $A\cup B$ contains those satisfying \emph{at least one}. Set difference $A\setminus B$ contains trials that are in $A$ but not $B$. Relative complement $A^c$ is meaningful only after choosing the universe $U$: here $A^c=U\setminus A$. We must never conflate the universe of trials in one session with all possible trials.

De Morgan's identities are
\[(A\cup B)^c=A^c\cap B^c,\qquad(A\cap B)^c=A^c\cup B^c.\]
For $U=\{1,2,3,4,5\}$, $A=\{1,3,5\}$, and $B=\{3,4\}$, we have $A\cap B=\{3\}$ and $A\cup B=\{1,3,4,5\}$. Work through the identities element by element before trusting the symbols.

\section{Families, indexing, and power sets}
A \emph{family} $(A_i)_{i\in I}$ is an indexed collection of sets. The union $\bigcup_{i\in I}A_i$ collects elements present in at least one $A_i$; the intersection $\bigcap_{i\in I}A_i$ collects elements present in every $A_i$. The \emph{power set} $\mathcal P(A)$ is the set of all subsets of $A$. For $A=\{u,v\}$ it is $\{\varnothing,\{u\},\{v\},\{u,v\}\}$. If $A$ has $n$ elements, its power set contains $2^n$ elements: each original element can independently be included or excluded.

\section{The significance for scientific data}
Many preprocessing operations change the set of data included in analysis. If $U$ is the set of detected units, then $U_{\rm quality}$ might retain only units meeting refractory-violation and signal-to-noise thresholds. Estimating a correlation on $U_{\rm quality}$ is not the same question as estimating it on $U$. Record the inclusion criteria as mathematical predicates; this makes sampling assumptions auditable.

\begin{workedexample}[Counting selections]
Suppose four electrodes pass quality control. A subset represents the electrodes chosen for a multichannel analysis; there are $2^4=16$ possible subsets, including the empty selection. Requiring at least one electrode yields 15. Requiring exactly two yields $\binom42=6$. This distinction prepares us for combinatorics and multiple-comparison correction.
\end{workedexample}
% --- expanded edition (October 2026) ---
\section{Algebra of sets}
The operations of the previous sections obey identities that parallel the logical equivalences of Chapter~\ref{ch:logic}. This parallel is not accidental: $x\in A\cap B$ means $x\in A\land x\in B$, $x\in A\cup B$ means $x\in A\lor x\in B$, and $x\in A^c$ means $\neg(x\in A)$. Every set identity can therefore be proved by translating to a logical equivalence about membership.

\begin{proposition}[Laws of the algebra of sets]
For subsets $A,B,C$ of a universe $U$:
\begin{enumerate}[leftmargin=2em]
\item commutativity: $A\cup B=B\cup A$ and $A\cap B=B\cap A$;
\item associativity: $(A\cup B)\cup C=A\cup(B\cup C)$, and similarly for $\cap$;
\item distributivity: $A\cap(B\cup C)=(A\cap B)\cup(A\cap C)$ and $A\cup(B\cap C)=(A\cup B)\cap(A\cup C)$;
\item identity and absorption: $A\cup\varnothing=A$, $A\cap U=A$, $A\cup(A\cap B)=A$;
\item complement: $A\cup A^c=U$, $A\cap A^c=\varnothing$, $(A^c)^c=A$.
\end{enumerate}
\end{proposition}
\begin{proof}
We prove the second distributive law; the others follow the same pattern. Let $x\in U$. Then $x\in A\cup(B\cap C)$ iff $x\in A\lor(x\in B\land x\in C)$. By the logical distributive law this is equivalent to $(x\in A\lor x\in B)\land(x\in A\lor x\in C)$, i.e.\ $x\in(A\cup B)\cap(A\cup C)$. Since the two sets have the same members, they are equal.
\end{proof}

\begin{proposition}[De Morgan for families]
For any indexed family $(A_i)_{i\in I}$ of subsets of $U$,
\[
\Bigl(\bigcup_{i\in I}A_i\Bigr)^c=\bigcap_{i\in I}A_i^c,\qquad
\Bigl(\bigcap_{i\in I}A_i\Bigr)^c=\bigcup_{i\in I}A_i^c.
\]
\end{proposition}
\begin{proof}
For the first identity, $x\notin\bigcup_iA_i$ iff $\neg\exists i\,(x\in A_i)$ iff $\forall i\,(x\notin A_i)$ iff $x\in\bigcap_iA_i^c$. The second follows by applying the first to the family $(A_i^c)$ and taking complements.
\end{proof}

\section{Symmetric difference and indicator functions}
The \emph{symmetric difference} $A\,\triangle\,B=(A\setminus B)\cup(B\setminus A)$ contains the elements in exactly one of the two sets. It measures disagreement between two classifications: if $A$ and $B$ are the spikes assigned to a unit by two sorting algorithms, $|A\triangle B|$ counts the spikes on which they disagree.

The \emph{indicator function} of $A\subseteq U$ is $\mathbf 1_A:U\to\{0,1\}$, with $\mathbf 1_A(x)=1$ if $x\in A$ and $0$ otherwise. Indicators turn set operations into arithmetic:
\[
\mathbf 1_{A\cap B}=\mathbf 1_A\mathbf 1_B,\qquad
\mathbf 1_{A^c}=1-\mathbf 1_A,\qquad
\mathbf 1_{A\cup B}=\mathbf 1_A+\mathbf 1_B-\mathbf 1_A\mathbf 1_B,\qquad
\mathbf 1_{A\triangle B}=|\mathbf 1_A-\mathbf 1_B|.
\]
Summing the union identity over a finite universe gives the \emph{inclusion--exclusion principle} $|A\cup B|=|A|+|B|-|A\cap B|$. Logical masks in numerical software are exactly indicator functions stored as arrays; an element-wise product of two masks is an intersection.

\section{Partitions}
A \emph{partition} of a set $U$ is a family of nonempty, pairwise disjoint subsets whose union is $U$. Trials partitioned by stimulus condition, spikes partitioned by unit identity, and time partitioned into bins are all partitions. If $\{B_1,\ldots,B_k\}$ partitions a finite $U$, then for any $A\subseteq U$,
\[
|A|=\sum_{j=1}^k|A\cap B_j|,
\]
because the sets $A\cap B_j$ are disjoint and their union is $A$. This additivity, generalized from counting to probability and integration, is the root of the law of total probability (Chapter~\ref{ch:probability}).

\section{Russell's paradox in detail}
Suppose there were a set $R$ whose members are exactly the sets that are not members of themselves: $x\in R\Leftrightarrow x\notin x$. Taking $x=R$ yields $R\in R\Leftrightarrow R\notin R$, a contradiction. Hence no such set exists. The \emph{axiom of separation} avoids the paradox by allowing $\{x\in A:P(x)\}$ only for an already given set $A$. Applied to any set $A$, the same argument shows that $\{x\in A:x\notin x\}$ is not a member of $A$; consequently, there is no set of all sets. In applied work this matters less as a danger than as a habit: every set in an analysis should be carved out of a specified universe.

\begin{workedexample}[Masks and exclusion rules]
A session has 400 trials. Let $M$ be trials with movement artifacts (37 trials), $L$ trials with lick contamination (52 trials), and suppose 15 trials have both. The excluded set is $M\cup L$, with $|M\cup L|=37+52-15=74$ trials, leaving $|(M\cup L)^c|=326$. Reporting ``37 and 52 trials were excluded'' would suggest 89 exclusions; inclusion--exclusion gives the correct count.
\end{workedexample}
\begin{exercises}
\item Prove $A\cap(B\cup C)=(A\cap B)\cup(A\cap C)$ from element membership.
\item Determine whether $\varnothing\in\{\varnothing\}$ and whether $\varnothing\subseteq\{\varnothing\}$.
\item Describe the difference between a population of neurons and its sampled subset.
\item Given three data-exclusion rules, express the accepted trials with set operations.
\item Prove $A\setminus(B\cup C)=(A\setminus B)\cap(A\setminus C)$.
\item Show that $A\triangle B=\varnothing$ if and only if $A=B$, and that $\triangle$ is associative.
\item Derive inclusion--exclusion for three sets from the indicator identity.
\item Show that $\{x\in A:x\notin x\}\notin A$ for every set $A$.
\item Two sorters assign 1200 and 1100 spikes to a unit and agree on 1000. Compute $|A\triangle B|$ and a normalized disagreement $|A\triangle B|/|A\cup B|$.
\end{exercises}
